% ============================================================================
% numbers.tex — EVERY quantitative value used in the manuscript, in one place.
%
% WHY THIS FILE EXISTS
%   The prose in main.tex contains no literal number. Every quantity is a macro
%   defined here (or in the machine-generated block this file pulls in), so a
%   value can be corrected in exactly one place and the manuscript follows.
%
% TWO-FILE STRUCTURE  (introduced 2026-08-08, coder's re-run merge)
%   numbers_reruns.tex   MACHINE-GENERATED, DO NOT HAND-EDIT. Every model-derived
%                        number: seed-averaged point estimate \X, inference-seed
%                        SD \XSD, bootstrap-over-subjects interval \XCI. Emitted by
%                        scripts/paper_reruns/emit_numbers.py from summary_ci.csv
%                        (336 rows) and contrast_tests.csv (186 rows). Every macro
%                        carries its source row as a preceding comment. To refresh:
%                        re-run the emitter and copy the file over.
%   numbers.tex (here)   Everything the emitter cannot produce: cohort counts from
%                        the manifests, chance baselines, literature values, the
%                        pre-correction leakage arm, and the prose-formatted
%                        wrappers around emitted intervals.
%
%   ORDER MATTERS. The \input comes first; this file may then \renewcommand an
%   emitted macro for presentation, but must never silently change its VALUE.
%   Any \renewcommand below that alters a number rather than its formatting is a
%   bug — the emitted block is the source of record for model numbers.
%
% WHAT THE INTERVALS MEAN
%   \XSD is spread across five inference seeds with the model WEIGHTS FROZEN —
%   it is the farthest-point-sampling draw, nothing else. \XCI is a bootstrap
%   over test SUBJECTS. The two answer different questions and are never pooled.
%   NEITHER captures training-seed variance: nothing was retrained. Methods says
%   so explicitly and must keep saying so.
%
% MARKERS (visible while \draftnumstrue, silent when set false)
%   \prov{x}   provisional — single unseeded run, no CI                (dagger)
%   \ver{x}    awaiting analyst verification against the manifests     (star)
%   \conflict{x} two sources disagree; must not ship until resolved    (ddagger)
%   \unsup{x}  NO EXPERIMENT BEHIND THIS NUMBER — cut it or run it     (section)
%   Drop the wrapper once a value is settled. The markers only mean something if
%   they are maintained honestly.
%
% PROVENANCE for the hand-maintained values below: dataset manifests (analyst,
% pending), EXPERIMENTAL_FINDINGS.md (read its 2026-08-04 editor's note first —
% its body still contains one withdrawn claim and one mislabelled experiment),
% results_analysis.md, refs.bib.
% ============================================================================

\input{numbers_reruns}
%   numbers_xsite.tex    MACHINE-GENERATED, DO NOT HAND-EDIT. The cross-site
%                        leave-one-site-out block. Emitted by
%                        scripts/paper_reruns/emit_numbers_xsite.py from
%                        summary_ci_xsite.csv (72 rows). Same \X / \XSD / \XCI
%                        convention. Group 1 (\xsGone*) and Group 2 (\xsGtwo*)
%                        are NEVER pooled — there is deliberately no macro
%                        spanning both. To refresh: re-run the emitter, copy over.
\input{numbers_xsite}
%   numbers_xsite_bin.tex  MACHINE-GENERATED, DO NOT HAND-EDIT. The SAME six
%                        leave-one-site-out folds run on the BINARY screening
%                        task (A2, affected vs unaffected within the cleft
%                        cohort). Emitted by the same script with --binary from
%                        summary_ci_xsite_bin.csv (72 rows). Its macros live in
%                        a separate \xsb* namespace precisely so that this file
%                        and numbers_xsite.tex can be \input together without
%                        either silently redefining the other's \xs<Site><Enc>
%                        names — the two blocks are DIFFERENT TASKS on the same
%                        folds, and a four-class value must never be readable
%                        through a binary macro name or vice versa.
%                        The within-site reference for this block is \ofcBin*
%                        (task A2), NOT \combAuc — the combined-cohort detection
%                        task was not run cross-site.
\input{numbers_xsite_bin}
%   numbers_xsitefix.tex / numbers_xsitefix_bin.tex
%                        MACHINE-GENERATED, DO NOT HAND-EDIT. Added 2026-09-08,
%                        thread facebase3d-xsitediag-2026-09-07. The SAME six
%                        folds, encoders and label collapses as the two blocks
%                        above, retrained with early stopping DISABLED after the
%                        published folds were found to have stopped at 21--46
%                        epochs, inside the ln(K) plateau. Emitted by the same
%                        script with --fix (and --fix --binary).
%                        THEY REBIND NOTHING. The corrected arm lives in \xsf*
%                        and \xsbf*; the published patience-20 arm keeps \xs*
%                        and \xsb*. Sec 2.7 quotes the corrected arm and cites
%                        the published one as the truncated comparison, so both
%                        have to remain readable and separately named.
%                        Each cell is still ONE training run, and the corrected
%                        runs are seeded while the published ones were not, so
%                        the difference between the two arms confounds the
%                        stopping rule with retraining variance. Methods says so.
% ==========================================================================
% numbers_from_xsitefix.tex — MACHINE-GENERATED. Do not hand-edit.
%
%   generator : scripts/paper_reruns/emit_numbers_xsite.py
%   source    : results/paper_reruns/summary_ci_xsitefix.csv  (72 rows)
%
% Leave-one-site-out cross-site transfer. TEST is one held-out
% FaceBase acquisition site in full; train/val are a stratified
% split over the remaining ten sites. The five sites failing the
% post-attrition per-class viability gate (CL, CO, NG, SF, TX) stay
% in the TRAINING pool and are never scored as folds.
%
% THERE WAS NO PRIOR RESULT HERE. The 'AUC 0.48-0.51' figures in
% earlier drafts came from the A1_S1/S2/S3 runs, which are
% class-imbalance variants on the full-cohort split, not site
% splits. \xsiteAucLo / \xsiteAucHi are NOT rebound by this block:
% new measurement, new names.
%
% THE TWO FOLD GROUPS ARE NEVER POOLED:
%   \xsGone*  PH, FC, PR — both datasets remain in training, so the
%             fold isolates SITE from DATASET
%   \xsGtwo*  GW, PT, LC — single-dataset; site and dataset move
%             together and are measured jointly
% No macro here spans both groups.
%
% READ ACCURACY AGAINST \xs<Site>Maj. Held-out sites are 72-92%
% Unaffected and prevalence differs per site, so accuracy is not
% comparable across folds. macro-AUC is the headline metric;
% its chance level is 0.50 regardless of prevalence.
%
% \X point estimate / \XSD inference-seed SD / \XCI bootstrap over
% test SUBJECTS. Never pooled.
%
% NOT MEASURED: training-seed variance. Each cell is ONE training
% run. Every interval is inference variance under frozen weights.
%
% THIS IS THE CORRECTED ARM (2026-09-07). Same six folds, same
% manifests, same encoders, same label collapse. TWO protocol
% changes relative to the \xs* block:
%   1. early stopping disabled (patience == epochs == 200).
%      The published folds stopped at 21-46 epochs, inside the
%      ln(K) plateau, so they measured the stopping rule as
%      much as site transfer.
%   2. training is SEEDED (--train-seed 1); the published runs
%      were not, and seeding also gives each dataloader worker
%      its own augmentation stream.
% Both arms are therefore ONE draw each. The difference between
% them is NOT separated into stopping rule vs retraining
% variance, and no macro here should be quoted as if it were.
% What licenses reading it as the stopping rule is the pattern,
% not any single cell: the gain is consistent across folds for
% the encoders that plateau, and absent for GeomMLP, which has
% no plateau and is the control.
% ==========================================================================

% --------------------------------------------------------------------------
% Fold sizes and majority-class baselines (post-attrition)
% --------------------------------------------------------------------------
% site PH | group 1 (site isolated) | datasets FB-56+FB-5A
\newcommand{\xsfPHNtest}{1960}
\newcommand{\xsfPHNtrain}{4178}
\newcommand{\xsfPHMaj}{0.800}  % majority-class rate = accuracy of a constant predictor

% site FC | group 1 (site isolated) | datasets FB-56+FB-5A
\newcommand{\xsfFCNtest}{1259}
\newcommand{\xsfFCNtrain}{4774}
\newcommand{\xsfFCMaj}{0.919}  % majority-class rate = accuracy of a constant predictor

% site PR | group 1 (site isolated) | datasets FB-56+FB-5A
\newcommand{\xsfPRNtest}{589}
\newcommand{\xsfPRNtrain}{5342}
\newcommand{\xsfPRMaj}{0.778}  % majority-class rate = accuracy of a constant predictor

% site GW | group 2 (site+dataset) | datasets FB-5A
\newcommand{\xsfGWNtest}{744}
\newcommand{\xsfGWNtrain}{5211}
\newcommand{\xsfGWMaj}{0.812}  % majority-class rate = accuracy of a constant predictor

% site PT | group 2 (site+dataset) | datasets FB-5A
\newcommand{\xsfPTNtest}{473}
\newcommand{\xsfPTNtrain}{5442}
\newcommand{\xsfPTMaj}{0.732}  % majority-class rate = accuracy of a constant predictor

% site LC | group 2 (site+dataset) | datasets FB-56
\newcommand{\xsfLCNtest}{120}
\newcommand{\xsfLCNtrain}{5742}
\newcommand{\xsfLCMaj}{0.708}  % majority-class rate = accuracy of a constant predictor

% --------------------------------------------------------------------------
% GROUP 1 — site isolated from dataset (both FB-56 and FB-5A remain in training)
% --------------------------------------------------------------------------
% summary_ci_xsitefix.csv | task=XS_fix_PH encoder=geommlp head=softmax level=scan metric=macro_auc n_seeds=1 n_units=1960 n_subjects=1959
\newcommand{\xsfPHGmlpAuc}{0.499}
\newcommand{\xsfPHGmlpAucSD}{0}  % exact 0: deterministic encoder, evaluated once
\newcommand{\xsfPHGmlpAucCI}{[0.468, 0.528]}  % 95% bootstrap over held-out-site test subjects

% summary_ci_xsitefix.csv | task=XS_fix_PH encoder=geommlp head=softmax level=scan metric=macro_f1 n_seeds=1 n_units=1960 n_subjects=1959
\newcommand{\xsfPHGmlpFone}{0.215}
\newcommand{\xsfPHGmlpFoneSD}{0}  % exact 0: deterministic encoder, evaluated once
\newcommand{\xsfPHGmlpFoneCI}{[0.196, 0.236]}  % 95% bootstrap over held-out-site test subjects

% summary_ci_xsitefix.csv | task=XS_fix_PH encoder=geommlp head=softmax level=scan metric=accuracy n_seeds=1 n_units=1960 n_subjects=1959
\newcommand{\xsfPHGmlpAcc}{0.384}
\newcommand{\xsfPHGmlpAccSD}{0}  % exact 0: deterministic encoder, evaluated once
\newcommand{\xsfPHGmlpAccCI}{[0.362, 0.406]}  % 95% bootstrap over held-out-site test subjects

% summary_ci_xsitefix.csv | task=XS_fix_PH encoder=pointnet head=softmax level=scan metric=macro_auc n_seeds=1 n_units=1960 n_subjects=1959
\newcommand{\xsfPHPnetAuc}{0.555}
\newcommand{\xsfPHPnetAucSD}{0}  % exact 0: deterministic encoder, evaluated once
\newcommand{\xsfPHPnetAucCI}{[0.524, 0.587]}  % 95% bootstrap over held-out-site test subjects

% summary_ci_xsitefix.csv | task=XS_fix_PH encoder=pointnet head=softmax level=scan metric=macro_f1 n_seeds=1 n_units=1960 n_subjects=1959
\newcommand{\xsfPHPnetFone}{0.273}
\newcommand{\xsfPHPnetFoneSD}{0}  % exact 0: deterministic encoder, evaluated once
\newcommand{\xsfPHPnetFoneCI}{[0.256, 0.292]}  % 95% bootstrap over held-out-site test subjects

% summary_ci_xsitefix.csv | task=XS_fix_PH encoder=pointnet head=softmax level=scan metric=accuracy n_seeds=1 n_units=1960 n_subjects=1959
\newcommand{\xsfPHPnetAcc}{0.586}
\newcommand{\xsfPHPnetAccSD}{0}  % exact 0: deterministic encoder, evaluated once
\newcommand{\xsfPHPnetAccCI}{[0.564, 0.607]}  % 95% bootstrap over held-out-site test subjects

% summary_ci_xsitefix.csv | task=XS_fix_PH encoder=dgcnn head=softmax level=scan metric=macro_auc n_seeds=1 n_units=1960 n_subjects=1959
\newcommand{\xsfPHDgcnnAuc}{0.614}
\newcommand{\xsfPHDgcnnAucSD}{0}  % exact 0: deterministic encoder, evaluated once
\newcommand{\xsfPHDgcnnAucCI}{[0.585, 0.644]}  % 95% bootstrap over held-out-site test subjects

% summary_ci_xsitefix.csv | task=XS_fix_PH encoder=dgcnn head=softmax level=scan metric=macro_f1 n_seeds=1 n_units=1960 n_subjects=1959
\newcommand{\xsfPHDgcnnFone}{0.263}
\newcommand{\xsfPHDgcnnFoneSD}{0}  % exact 0: deterministic encoder, evaluated once
\newcommand{\xsfPHDgcnnFoneCI}{[0.247, 0.280]}  % 95% bootstrap over held-out-site test subjects

% summary_ci_xsitefix.csv | task=XS_fix_PH encoder=dgcnn head=softmax level=scan metric=accuracy n_seeds=1 n_units=1960 n_subjects=1959
\newcommand{\xsfPHDgcnnAcc}{0.585}
\newcommand{\xsfPHDgcnnAccSD}{0}  % exact 0: deterministic encoder, evaluated once
\newcommand{\xsfPHDgcnnAccCI}{[0.564, 0.607]}  % 95% bootstrap over held-out-site test subjects

% summary_ci_xsitefix.csv | task=XS_fix_PH encoder=pointnet2 head=softmax level=scan metric=macro_auc n_seeds=5 n_units=1960 n_subjects=1959
\newcommand{\xsfPHPtwoAuc}{0.639}
\newcommand{\xsfPHPtwoAucSD}{0.0027}  % inference-seed SD, weights frozen
\newcommand{\xsfPHPtwoAucCI}{[0.610, 0.669]}  % 95% bootstrap over held-out-site test subjects

% summary_ci_xsitefix.csv | task=XS_fix_PH encoder=pointnet2 head=softmax level=scan metric=macro_f1 n_seeds=5 n_units=1960 n_subjects=1959
\newcommand{\xsfPHPtwoFone}{0.280}
\newcommand{\xsfPHPtwoFoneSD}{0.0061}  % inference-seed SD, weights frozen
\newcommand{\xsfPHPtwoFoneCI}{[0.261, 0.300]}  % 95% bootstrap over held-out-site test subjects

% summary_ci_xsitefix.csv | task=XS_fix_PH encoder=pointnet2 head=softmax level=scan metric=accuracy n_seeds=5 n_units=1960 n_subjects=1959
\newcommand{\xsfPHPtwoAcc}{0.479}
\newcommand{\xsfPHPtwoAccSD}{0.0062}  % inference-seed SD, weights frozen
\newcommand{\xsfPHPtwoAccCI}{[0.460, 0.498]}  % 95% bootstrap over held-out-site test subjects

% summary_ci_xsitefix.csv | task=XS_fix_FC encoder=geommlp head=softmax level=scan metric=macro_auc n_seeds=1 n_units=1259 n_subjects=1259
\newcommand{\xsfFCGmlpAuc}{0.617}
\newcommand{\xsfFCGmlpAucSD}{0}  % exact 0: deterministic encoder, evaluated once
\newcommand{\xsfFCGmlpAucCI}{[0.563, 0.678]}  % 95% bootstrap over held-out-site test subjects

% summary_ci_xsitefix.csv | task=XS_fix_FC encoder=geommlp head=softmax level=scan metric=macro_f1 n_seeds=1 n_units=1259 n_subjects=1259
\newcommand{\xsfFCGmlpFone}{0.212}
\newcommand{\xsfFCGmlpFoneSD}{0}  % exact 0: deterministic encoder, evaluated once
\newcommand{\xsfFCGmlpFoneCI}{[0.193, 0.233]}  % 95% bootstrap over held-out-site test subjects

% summary_ci_xsitefix.csv | task=XS_fix_FC encoder=geommlp head=softmax level=scan metric=accuracy n_seeds=1 n_units=1259 n_subjects=1259
\newcommand{\xsfFCGmlpAcc}{0.531}
\newcommand{\xsfFCGmlpAccSD}{0}  % exact 0: deterministic encoder, evaluated once
\newcommand{\xsfFCGmlpAccCI}{[0.503, 0.558]}  % 95% bootstrap over held-out-site test subjects

% summary_ci_xsitefix.csv | task=XS_fix_FC encoder=pointnet head=softmax level=scan metric=macro_auc n_seeds=1 n_units=1259 n_subjects=1259
\newcommand{\xsfFCPnetAuc}{0.524}
\newcommand{\xsfFCPnetAucSD}{0}  % exact 0: deterministic encoder, evaluated once
\newcommand{\xsfFCPnetAucCI}{[0.473, 0.577]}  % 95% bootstrap over held-out-site test subjects

% summary_ci_xsitefix.csv | task=XS_fix_FC encoder=pointnet head=softmax level=scan metric=macro_f1 n_seeds=1 n_units=1259 n_subjects=1259
\newcommand{\xsfFCPnetFone}{0.217}
\newcommand{\xsfFCPnetFoneSD}{0}  % exact 0: deterministic encoder, evaluated once
\newcommand{\xsfFCPnetFoneCI}{[0.207, 0.228]}  % 95% bootstrap over held-out-site test subjects

% summary_ci_xsitefix.csv | task=XS_fix_FC encoder=pointnet head=softmax level=scan metric=accuracy n_seeds=1 n_units=1259 n_subjects=1259
\newcommand{\xsfFCPnetAcc}{0.630}
\newcommand{\xsfFCPnetAccSD}{0}  % exact 0: deterministic encoder, evaluated once
\newcommand{\xsfFCPnetAccCI}{[0.602, 0.655]}  % 95% bootstrap over held-out-site test subjects

% summary_ci_xsitefix.csv | task=XS_fix_FC encoder=dgcnn head=softmax level=scan metric=macro_auc n_seeds=1 n_units=1259 n_subjects=1259
\newcommand{\xsfFCDgcnnAuc}{0.492}
\newcommand{\xsfFCDgcnnAucSD}{0}  % exact 0: deterministic encoder, evaluated once
\newcommand{\xsfFCDgcnnAucCI}{[0.433, 0.549]}  % 95% bootstrap over held-out-site test subjects

% summary_ci_xsitefix.csv | task=XS_fix_FC encoder=dgcnn head=softmax level=scan metric=macro_f1 n_seeds=1 n_units=1259 n_subjects=1259
\newcommand{\xsfFCDgcnnFone}{0.216}
\newcommand{\xsfFCDgcnnFoneSD}{0}  % exact 0: deterministic encoder, evaluated once
\newcommand{\xsfFCDgcnnFoneCI}{[0.198, 0.235]}  % 95% bootstrap over held-out-site test subjects

% summary_ci_xsitefix.csv | task=XS_fix_FC encoder=dgcnn head=softmax level=scan metric=accuracy n_seeds=1 n_units=1259 n_subjects=1259
\newcommand{\xsfFCDgcnnAcc}{0.562}
\newcommand{\xsfFCDgcnnAccSD}{0}  % exact 0: deterministic encoder, evaluated once
\newcommand{\xsfFCDgcnnAccCI}{[0.534, 0.589]}  % 95% bootstrap over held-out-site test subjects

% summary_ci_xsitefix.csv | task=XS_fix_FC encoder=pointnet2 head=softmax level=scan metric=macro_auc n_seeds=5 n_units=1259 n_subjects=1259
\newcommand{\xsfFCPtwoAuc}{0.598}
\newcommand{\xsfFCPtwoAucSD}{0.0049}  % inference-seed SD, weights frozen
\newcommand{\xsfFCPtwoAucCI}{[0.536, 0.657]}  % 95% bootstrap over held-out-site test subjects

% summary_ci_xsitefix.csv | task=XS_fix_FC encoder=pointnet2 head=softmax level=scan metric=macro_f1 n_seeds=5 n_units=1259 n_subjects=1259
\newcommand{\xsfFCPtwoFone}{0.186}
\newcommand{\xsfFCPtwoFoneSD}{0.0022}  % inference-seed SD, weights frozen
\newcommand{\xsfFCPtwoFoneCI}{[0.176, 0.195]}  % 95% bootstrap over held-out-site test subjects

% summary_ci_xsitefix.csv | task=XS_fix_FC encoder=pointnet2 head=softmax level=scan metric=accuracy n_seeds=5 n_units=1259 n_subjects=1259
\newcommand{\xsfFCPtwoAcc}{0.480}
\newcommand{\xsfFCPtwoAccSD}{0.0072}  % inference-seed SD, weights frozen
\newcommand{\xsfFCPtwoAccCI}{[0.456, 0.504]}  % 95% bootstrap over held-out-site test subjects

% summary_ci_xsitefix.csv | task=XS_fix_PR encoder=geommlp head=softmax level=scan metric=macro_auc n_seeds=1 n_units=589 n_subjects=589
\newcommand{\xsfPRGmlpAuc}{0.621}
\newcommand{\xsfPRGmlpAucSD}{0}  % exact 0: deterministic encoder, evaluated once
\newcommand{\xsfPRGmlpAucCI}{[0.575, 0.671]}  % 95% bootstrap over held-out-site test subjects

% summary_ci_xsitefix.csv | task=XS_fix_PR encoder=geommlp head=softmax level=scan metric=macro_f1 n_seeds=1 n_units=589 n_subjects=589
\newcommand{\xsfPRGmlpFone}{0.305}
\newcommand{\xsfPRGmlpFoneSD}{0}  % exact 0: deterministic encoder, evaluated once
\newcommand{\xsfPRGmlpFoneCI}{[0.256, 0.357]}  % 95% bootstrap over held-out-site test subjects

% summary_ci_xsitefix.csv | task=XS_fix_PR encoder=geommlp head=softmax level=scan metric=accuracy n_seeds=1 n_units=589 n_subjects=589
\newcommand{\xsfPRGmlpAcc}{0.543}
\newcommand{\xsfPRGmlpAccSD}{0}  % exact 0: deterministic encoder, evaluated once
\newcommand{\xsfPRGmlpAccCI}{[0.504, 0.584]}  % 95% bootstrap over held-out-site test subjects

% summary_ci_xsitefix.csv | task=XS_fix_PR encoder=pointnet head=softmax level=scan metric=macro_auc n_seeds=1 n_units=589 n_subjects=589
\newcommand{\xsfPRPnetAuc}{0.576}
\newcommand{\xsfPRPnetAucSD}{0}  % exact 0: deterministic encoder, evaluated once
\newcommand{\xsfPRPnetAucCI}{[0.526, 0.630]}  % 95% bootstrap over held-out-site test subjects

% summary_ci_xsitefix.csv | task=XS_fix_PR encoder=pointnet head=softmax level=scan metric=macro_f1 n_seeds=1 n_units=589 n_subjects=589
\newcommand{\xsfPRPnetFone}{0.257}
\newcommand{\xsfPRPnetFoneSD}{0}  % exact 0: deterministic encoder, evaluated once
\newcommand{\xsfPRPnetFoneCI}{[0.228, 0.292]}  % 95% bootstrap over held-out-site test subjects

% summary_ci_xsitefix.csv | task=XS_fix_PR encoder=pointnet head=softmax level=scan metric=accuracy n_seeds=1 n_units=589 n_subjects=589
\newcommand{\xsfPRPnetAcc}{0.599}
\newcommand{\xsfPRPnetAccSD}{0}  % exact 0: deterministic encoder, evaluated once
\newcommand{\xsfPRPnetAccCI}{[0.560, 0.640]}  % 95% bootstrap over held-out-site test subjects

% summary_ci_xsitefix.csv | task=XS_fix_PR encoder=dgcnn head=softmax level=scan metric=macro_auc n_seeds=1 n_units=589 n_subjects=589
\newcommand{\xsfPRDgcnnAuc}{0.734}
\newcommand{\xsfPRDgcnnAucSD}{0}  % exact 0: deterministic encoder, evaluated once
\newcommand{\xsfPRDgcnnAucCI}{[0.693, 0.774]}  % 95% bootstrap over held-out-site test subjects

% summary_ci_xsitefix.csv | task=XS_fix_PR encoder=dgcnn head=softmax level=scan metric=macro_f1 n_seeds=1 n_units=589 n_subjects=589
\newcommand{\xsfPRDgcnnFone}{0.329}
\newcommand{\xsfPRDgcnnFoneSD}{0}  % exact 0: deterministic encoder, evaluated once
\newcommand{\xsfPRDgcnnFoneCI}{[0.288, 0.375]}  % 95% bootstrap over held-out-site test subjects

% summary_ci_xsitefix.csv | task=XS_fix_PR encoder=dgcnn head=softmax level=scan metric=accuracy n_seeds=1 n_units=589 n_subjects=589
\newcommand{\xsfPRDgcnnAcc}{0.699}
\newcommand{\xsfPRDgcnnAccSD}{0}  % exact 0: deterministic encoder, evaluated once
\newcommand{\xsfPRDgcnnAccCI}{[0.664, 0.737]}  % 95% bootstrap over held-out-site test subjects

% summary_ci_xsitefix.csv | task=XS_fix_PR encoder=pointnet2 head=softmax level=scan metric=macro_auc n_seeds=5 n_units=589 n_subjects=589
\newcommand{\xsfPRPtwoAuc}{0.704}
\newcommand{\xsfPRPtwoAucSD}{0.0068}  % inference-seed SD, weights frozen
\newcommand{\xsfPRPtwoAucCI}{[0.660, 0.747]}  % 95% bootstrap over held-out-site test subjects

% summary_ci_xsitefix.csv | task=XS_fix_PR encoder=pointnet2 head=softmax level=scan metric=macro_f1 n_seeds=5 n_units=589 n_subjects=589
\newcommand{\xsfPRPtwoFone}{0.321}
\newcommand{\xsfPRPtwoFoneSD}{0.0085}  % inference-seed SD, weights frozen
\newcommand{\xsfPRPtwoFoneCI}{[0.291, 0.353]}  % 95% bootstrap over held-out-site test subjects

% summary_ci_xsitefix.csv | task=XS_fix_PR encoder=pointnet2 head=softmax level=scan metric=accuracy n_seeds=5 n_units=589 n_subjects=589
\newcommand{\xsfPRPtwoAcc}{0.626}
\newcommand{\xsfPRPtwoAccSD}{0.0078}  % inference-seed SD, weights frozen
\newcommand{\xsfPRPtwoAccCI}{[0.591, 0.661]}  % 95% bootstrap over held-out-site test subjects

% --------------------------------------------------------------------------
% GROUP 2 — site and dataset confounded (single-dataset sites)
% --------------------------------------------------------------------------
% summary_ci_xsitefix.csv | task=XS_fix_GW encoder=geommlp head=softmax level=scan metric=macro_auc n_seeds=1 n_units=744 n_subjects=744
\newcommand{\xsfGWGmlpAuc}{0.517}
\newcommand{\xsfGWGmlpAucSD}{0}  % exact 0: deterministic encoder, evaluated once
\newcommand{\xsfGWGmlpAucCI}{[0.474, 0.560]}  % 95% bootstrap over held-out-site test subjects

% summary_ci_xsitefix.csv | task=XS_fix_GW encoder=geommlp head=softmax level=scan metric=macro_f1 n_seeds=1 n_units=744 n_subjects=744
\newcommand{\xsfGWGmlpFone}{0.204}
\newcommand{\xsfGWGmlpFoneSD}{0}  % exact 0: deterministic encoder, evaluated once
\newcommand{\xsfGWGmlpFoneCI}{[0.178, 0.231]}  % 95% bootstrap over held-out-site test subjects

% summary_ci_xsitefix.csv | task=XS_fix_GW encoder=geommlp head=softmax level=scan metric=accuracy n_seeds=1 n_units=744 n_subjects=744
\newcommand{\xsfGWGmlpAcc}{0.370}
\newcommand{\xsfGWGmlpAccSD}{0}  % exact 0: deterministic encoder, evaluated once
\newcommand{\xsfGWGmlpAccCI}{[0.336, 0.403]}  % 95% bootstrap over held-out-site test subjects

% summary_ci_xsitefix.csv | task=XS_fix_GW encoder=pointnet head=softmax level=scan metric=macro_auc n_seeds=1 n_units=744 n_subjects=744
\newcommand{\xsfGWPnetAuc}{0.510}
\newcommand{\xsfGWPnetAucSD}{0}  % exact 0: deterministic encoder, evaluated once
\newcommand{\xsfGWPnetAucCI}{[0.463, 0.556]}  % 95% bootstrap over held-out-site test subjects

% summary_ci_xsitefix.csv | task=XS_fix_GW encoder=pointnet head=softmax level=scan metric=macro_f1 n_seeds=1 n_units=744 n_subjects=744
\newcommand{\xsfGWPnetFone}{0.244}
\newcommand{\xsfGWPnetFoneSD}{0}  % exact 0: deterministic encoder, evaluated once
\newcommand{\xsfGWPnetFoneCI}{[0.226, 0.263]}  % 95% bootstrap over held-out-site test subjects

% summary_ci_xsitefix.csv | task=XS_fix_GW encoder=pointnet head=softmax level=scan metric=accuracy n_seeds=1 n_units=744 n_subjects=744
\newcommand{\xsfGWPnetAcc}{0.706}
\newcommand{\xsfGWPnetAccSD}{0}  % exact 0: deterministic encoder, evaluated once
\newcommand{\xsfGWPnetAccCI}{[0.672, 0.738]}  % 95% bootstrap over held-out-site test subjects

% summary_ci_xsitefix.csv | task=XS_fix_GW encoder=dgcnn head=softmax level=scan metric=macro_auc n_seeds=1 n_units=744 n_subjects=744
\newcommand{\xsfGWDgcnnAuc}{0.674}
\newcommand{\xsfGWDgcnnAucSD}{0}  % exact 0: deterministic encoder, evaluated once
\newcommand{\xsfGWDgcnnAucCI}{[0.633, 0.714]}  % 95% bootstrap over held-out-site test subjects

% summary_ci_xsitefix.csv | task=XS_fix_GW encoder=dgcnn head=softmax level=scan metric=macro_f1 n_seeds=1 n_units=744 n_subjects=744
\newcommand{\xsfGWDgcnnFone}{0.331}
\newcommand{\xsfGWDgcnnFoneSD}{0}  % exact 0: deterministic encoder, evaluated once
\newcommand{\xsfGWDgcnnFoneCI}{[0.291, 0.372]}  % 95% bootstrap over held-out-site test subjects

% summary_ci_xsitefix.csv | task=XS_fix_GW encoder=dgcnn head=softmax level=scan metric=accuracy n_seeds=1 n_units=744 n_subjects=744
\newcommand{\xsfGWDgcnnAcc}{0.638}
\newcommand{\xsfGWDgcnnAccSD}{0}  % exact 0: deterministic encoder, evaluated once
\newcommand{\xsfGWDgcnnAccCI}{[0.603, 0.673]}  % 95% bootstrap over held-out-site test subjects

% summary_ci_xsitefix.csv | task=XS_fix_GW encoder=pointnet2 head=softmax level=scan metric=macro_auc n_seeds=5 n_units=744 n_subjects=744
\newcommand{\xsfGWPtwoAuc}{0.639}
\newcommand{\xsfGWPtwoAucSD}{0.0107}  % inference-seed SD, weights frozen
\newcommand{\xsfGWPtwoAucCI}{[0.594, 0.682]}  % 95% bootstrap over held-out-site test subjects

% summary_ci_xsitefix.csv | task=XS_fix_GW encoder=pointnet2 head=softmax level=scan metric=macro_f1 n_seeds=5 n_units=744 n_subjects=744
\newcommand{\xsfGWPtwoFone}{0.279}
\newcommand{\xsfGWPtwoFoneSD}{0.0107}  % inference-seed SD, weights frozen
\newcommand{\xsfGWPtwoFoneCI}{[0.252, 0.307]}  % 95% bootstrap over held-out-site test subjects

% summary_ci_xsitefix.csv | task=XS_fix_GW encoder=pointnet2 head=softmax level=scan metric=accuracy n_seeds=5 n_units=744 n_subjects=744
\newcommand{\xsfGWPtwoAcc}{0.516}
\newcommand{\xsfGWPtwoAccSD}{0.0128}  % inference-seed SD, weights frozen
\newcommand{\xsfGWPtwoAccCI}{[0.485, 0.547]}  % 95% bootstrap over held-out-site test subjects

% summary_ci_xsitefix.csv | task=XS_fix_PT encoder=geommlp head=softmax level=scan metric=macro_auc n_seeds=1 n_units=473 n_subjects=473
\newcommand{\xsfPTGmlpAuc}{0.568}
\newcommand{\xsfPTGmlpAucSD}{0}  % exact 0: deterministic encoder, evaluated once
\newcommand{\xsfPTGmlpAucCI}{[0.509, 0.628]}  % 95% bootstrap over held-out-site test subjects

% summary_ci_xsitefix.csv | task=XS_fix_PT encoder=geommlp head=softmax level=scan metric=macro_f1 n_seeds=1 n_units=473 n_subjects=473
\newcommand{\xsfPTGmlpFone}{0.264}
\newcommand{\xsfPTGmlpFoneSD}{0}  % exact 0: deterministic encoder, evaluated once
\newcommand{\xsfPTGmlpFoneCI}{[0.233, 0.297]}  % 95% bootstrap over held-out-site test subjects

% summary_ci_xsitefix.csv | task=XS_fix_PT encoder=geommlp head=softmax level=scan metric=accuracy n_seeds=1 n_units=473 n_subjects=473
\newcommand{\xsfPTGmlpAcc}{0.588}
\newcommand{\xsfPTGmlpAccSD}{0}  % exact 0: deterministic encoder, evaluated once
\newcommand{\xsfPTGmlpAccCI}{[0.543, 0.632]}  % 95% bootstrap over held-out-site test subjects

% summary_ci_xsitefix.csv | task=XS_fix_PT encoder=pointnet head=softmax level=scan metric=macro_auc n_seeds=1 n_units=473 n_subjects=473
\newcommand{\xsfPTPnetAuc}{0.510}
\newcommand{\xsfPTPnetAucSD}{0}  % exact 0: deterministic encoder, evaluated once
\newcommand{\xsfPTPnetAucCI}{[0.452, 0.563]}  % 95% bootstrap over held-out-site test subjects

% summary_ci_xsitefix.csv | task=XS_fix_PT encoder=pointnet head=softmax level=scan metric=macro_f1 n_seeds=1 n_units=473 n_subjects=473
\newcommand{\xsfPTPnetFone}{0.233}
\newcommand{\xsfPTPnetFoneSD}{0}  % exact 0: deterministic encoder, evaluated once
\newcommand{\xsfPTPnetFoneCI}{[0.207, 0.266]}  % 95% bootstrap over held-out-site test subjects

% summary_ci_xsitefix.csv | task=XS_fix_PT encoder=pointnet head=softmax level=scan metric=accuracy n_seeds=1 n_units=473 n_subjects=473
\newcommand{\xsfPTPnetAcc}{0.693}
\newcommand{\xsfPTPnetAccSD}{0}  % exact 0: deterministic encoder, evaluated once
\newcommand{\xsfPTPnetAccCI}{[0.649, 0.734]}  % 95% bootstrap over held-out-site test subjects

% summary_ci_xsitefix.csv | task=XS_fix_PT encoder=dgcnn head=softmax level=scan metric=macro_auc n_seeds=1 n_units=473 n_subjects=473
\newcommand{\xsfPTDgcnnAuc}{0.620}
\newcommand{\xsfPTDgcnnAucSD}{0}  % exact 0: deterministic encoder, evaluated once
\newcommand{\xsfPTDgcnnAucCI}{[0.563, 0.675]}  % 95% bootstrap over held-out-site test subjects

% summary_ci_xsitefix.csv | task=XS_fix_PT encoder=dgcnn head=softmax level=scan metric=macro_f1 n_seeds=1 n_units=473 n_subjects=473
\newcommand{\xsfPTDgcnnFone}{0.290}
\newcommand{\xsfPTDgcnnFoneSD}{0}  % exact 0: deterministic encoder, evaluated once
\newcommand{\xsfPTDgcnnFoneCI}{[0.257, 0.323]}  % 95% bootstrap over held-out-site test subjects

% summary_ci_xsitefix.csv | task=XS_fix_PT encoder=dgcnn head=softmax level=scan metric=accuracy n_seeds=1 n_units=473 n_subjects=473
\newcommand{\xsfPTDgcnnAcc}{0.636}
\newcommand{\xsfPTDgcnnAccSD}{0}  % exact 0: deterministic encoder, evaluated once
\newcommand{\xsfPTDgcnnAccCI}{[0.592, 0.677]}  % 95% bootstrap over held-out-site test subjects

% summary_ci_xsitefix.csv | task=XS_fix_PT encoder=pointnet2 head=softmax level=scan metric=macro_auc n_seeds=5 n_units=473 n_subjects=473
\newcommand{\xsfPTPtwoAuc}{0.602}
\newcommand{\xsfPTPtwoAucSD}{0.0050}  % inference-seed SD, weights frozen
\newcommand{\xsfPTPtwoAucCI}{[0.549, 0.655]}  % 95% bootstrap over held-out-site test subjects

% summary_ci_xsitefix.csv | task=XS_fix_PT encoder=pointnet2 head=softmax level=scan metric=macro_f1 n_seeds=5 n_units=473 n_subjects=473
\newcommand{\xsfPTPtwoFone}{0.328}
\newcommand{\xsfPTPtwoFoneSD}{0.0229}  % inference-seed SD, weights frozen
\newcommand{\xsfPTPtwoFoneCI}{[0.286, 0.378]}  % 95% bootstrap over held-out-site test subjects

% summary_ci_xsitefix.csv | task=XS_fix_PT encoder=pointnet2 head=softmax level=scan metric=accuracy n_seeds=5 n_units=473 n_subjects=473
\newcommand{\xsfPTPtwoAcc}{0.613}
\newcommand{\xsfPTPtwoAccSD}{0.0162}  % inference-seed SD, weights frozen
\newcommand{\xsfPTPtwoAccCI}{[0.576, 0.649]}  % 95% bootstrap over held-out-site test subjects

% summary_ci_xsitefix.csv | task=XS_fix_LC encoder=geommlp head=softmax level=scan metric=macro_auc n_seeds=1 n_units=120 n_subjects=120
\newcommand{\xsfLCGmlpAuc}{0.436}
\newcommand{\xsfLCGmlpAucSD}{0}  % exact 0: deterministic encoder, evaluated once
\newcommand{\xsfLCGmlpAucCI}{[0.344, 0.528]}  % 95% bootstrap over held-out-site test subjects

% summary_ci_xsitefix.csv | task=XS_fix_LC encoder=geommlp head=softmax level=scan metric=macro_f1 n_seeds=1 n_units=120 n_subjects=120
\newcommand{\xsfLCGmlpFone}{0.222}
\newcommand{\xsfLCGmlpFoneSD}{0}  % exact 0: deterministic encoder, evaluated once
\newcommand{\xsfLCGmlpFoneCI}{[0.189, 0.268]}  % 95% bootstrap over held-out-site test subjects

% summary_ci_xsitefix.csv | task=XS_fix_LC encoder=geommlp head=softmax level=scan metric=accuracy n_seeds=1 n_units=120 n_subjects=120
\newcommand{\xsfLCGmlpAcc}{0.667}
\newcommand{\xsfLCGmlpAccSD}{0}  % exact 0: deterministic encoder, evaluated once
\newcommand{\xsfLCGmlpAccCI}{[0.575, 0.750]}  % 95% bootstrap over held-out-site test subjects

% summary_ci_xsitefix.csv | task=XS_fix_LC encoder=pointnet head=softmax level=scan metric=macro_auc n_seeds=1 n_units=120 n_subjects=120
\newcommand{\xsfLCPnetAuc}{0.479}
\newcommand{\xsfLCPnetAucSD}{0}  % exact 0: deterministic encoder, evaluated once
\newcommand{\xsfLCPnetAucCI}{[0.370, 0.584]}  % 95% bootstrap over held-out-site test subjects

% summary_ci_xsitefix.csv | task=XS_fix_LC encoder=pointnet head=softmax level=scan metric=macro_f1 n_seeds=1 n_units=120 n_subjects=120
\newcommand{\xsfLCPnetFone}{0.247}
\newcommand{\xsfLCPnetFoneSD}{0}  % exact 0: deterministic encoder, evaluated once
\newcommand{\xsfLCPnetFoneCI}{[0.196, 0.307]}  % 95% bootstrap over held-out-site test subjects

% summary_ci_xsitefix.csv | task=XS_fix_LC encoder=pointnet head=softmax level=scan metric=accuracy n_seeds=1 n_units=120 n_subjects=120
\newcommand{\xsfLCPnetAcc}{0.683}
\newcommand{\xsfLCPnetAccSD}{0}  % exact 0: deterministic encoder, evaluated once
\newcommand{\xsfLCPnetAccCI}{[0.592, 0.767]}  % 95% bootstrap over held-out-site test subjects

% summary_ci_xsitefix.csv | task=XS_fix_LC encoder=dgcnn head=softmax level=scan metric=macro_auc n_seeds=1 n_units=120 n_subjects=120
\newcommand{\xsfLCDgcnnAuc}{0.516}
\newcommand{\xsfLCDgcnnAucSD}{0}  % exact 0: deterministic encoder, evaluated once
\newcommand{\xsfLCDgcnnAucCI}{[0.423, 0.609]}  % 95% bootstrap over held-out-site test subjects

% summary_ci_xsitefix.csv | task=XS_fix_LC encoder=dgcnn head=softmax level=scan metric=macro_f1 n_seeds=1 n_units=120 n_subjects=120
\newcommand{\xsfLCDgcnnFone}{0.237}
\newcommand{\xsfLCDgcnnFoneSD}{0}  % exact 0: deterministic encoder, evaluated once
\newcommand{\xsfLCDgcnnFoneCI}{[0.184, 0.288]}  % 95% bootstrap over held-out-site test subjects

% summary_ci_xsitefix.csv | task=XS_fix_LC encoder=dgcnn head=softmax level=scan metric=accuracy n_seeds=1 n_units=120 n_subjects=120
\newcommand{\xsfLCDgcnnAcc}{0.558}
\newcommand{\xsfLCDgcnnAccSD}{0}  % exact 0: deterministic encoder, evaluated once
\newcommand{\xsfLCDgcnnAccCI}{[0.467, 0.650]}  % 95% bootstrap over held-out-site test subjects

% summary_ci_xsitefix.csv | task=XS_fix_LC encoder=pointnet2 head=softmax level=scan metric=macro_auc n_seeds=5 n_units=120 n_subjects=120
\newcommand{\xsfLCPtwoAuc}{0.545}
\newcommand{\xsfLCPtwoAucSD}{0.0275}  % inference-seed SD, weights frozen
\newcommand{\xsfLCPtwoAucCI}{[0.460, 0.631]}  % 95% bootstrap over held-out-site test subjects

% summary_ci_xsitefix.csv | task=XS_fix_LC encoder=pointnet2 head=softmax level=scan metric=macro_f1 n_seeds=5 n_units=120 n_subjects=120
\newcommand{\xsfLCPtwoFone}{0.256}
\newcommand{\xsfLCPtwoFoneSD}{0.0154}  % inference-seed SD, weights frozen
\newcommand{\xsfLCPtwoFoneCI}{[0.202, 0.305]}  % 95% bootstrap over held-out-site test subjects

% summary_ci_xsitefix.csv | task=XS_fix_LC encoder=pointnet2 head=softmax level=scan metric=accuracy n_seeds=5 n_units=120 n_subjects=120
\newcommand{\xsfLCPtwoAcc}{0.642}
\newcommand{\xsfLCPtwoAccSD}{0.0177}  % inference-seed SD, weights frozen
\newcommand{\xsfLCPtwoAccCI}{[0.557, 0.720]}  % 95% bootstrap over held-out-site test subjects

% --------------------------------------------------------------------------
% Group-level macro-AUC ranges over folds, PointNet++ only.
% Ranges, not means: three folds is too few for an average to mean
% anything, and averaging would also hide a fold that behaves
% differently from the other two.
% --------------------------------------------------------------------------
% folds FC, PH, PR | encoder pointnet2 | metric macro_auc | level scan
\newcommand{\xsfGoneAucLo}{0.598}
\newcommand{\xsfGoneAucHi}{0.704}
\newcommand{\xsfGoneNfolds}{3}
\newcommand{\xsfGoneNaboveChance}{3}  % folds whose 95% CI lies entirely above 0.50

% folds GW, LC, PT | encoder pointnet2 | metric macro_auc | level scan
\newcommand{\xsfGtwoAucLo}{0.545}
\newcommand{\xsfGtwoAucHi}{0.639}
\newcommand{\xsfGtwoNfolds}{3}
\newcommand{\xsfGtwoNaboveChance}{2}  % folds whose 95% CI lies entirely above 0.50

% --------------------------------------------------------------------------
% Census over all fold x encoder cells, by their own bootstrap CI.
% --------------------------------------------------------------------------
\newcommand{\xsfNcells}{24}
\newcommand{\xsfNcellsAbove}{14}  % CI entirely above chance
\newcommand{\xsfNcellsSpan}{10}  % CI contains chance
\newcommand{\xsfNcellsBelow}{0}  % CI entirely below chance
\newcommand{\xsfGmlpNaboveChance}{3}  % of 6 folds, geommlp
\newcommand{\xsfGmlpAucLo}{0.436}
\newcommand{\xsfGmlpAucHi}{0.621}
\newcommand{\xsfPnetNaboveChance}{2}  % of 6 folds, pointnet
\newcommand{\xsfPnetAucLo}{0.479}
\newcommand{\xsfPnetAucHi}{0.576}
\newcommand{\xsfDgcnnNaboveChance}{4}  % of 6 folds, dgcnn
\newcommand{\xsfDgcnnAucLo}{0.492}
\newcommand{\xsfDgcnnAucHi}{0.734}
\newcommand{\xsfPtwoNaboveChance}{5}  % of 6 folds, pointnet2
\newcommand{\xsfPtwoAucLo}{0.545}
\newcommand{\xsfPtwoAucHi}{0.704}

% --------------------------------------------------------------------------
% Attenuation vs within-site (A1_fix, 5-retraining mean, scan level, softmax head).
% A fold MEAN is used only for this delta and never as a headline: the
% folds differ in size, prevalence and whether dataset moves with site,
% so they are not exchangeable and the prose reports ranges.
% --------------------------------------------------------------------------
% geommlp: within-site 0.590 -> fold mean 0.543
\newcommand{\xsfGmlpWithin}{0.590}
\newcommand{\xsfGmlpFoldMean}{0.543}
\newcommand{\xsfGmlpAtten}{4.7}  % percentage points, within-site minus fold mean
% pointnet: within-site 0.531 -> fold mean 0.526
\newcommand{\xsfPnetWithin}{0.531}
\newcommand{\xsfPnetFoldMean}{0.526}
\newcommand{\xsfPnetAtten}{0.5}  % percentage points, within-site minus fold mean
% dgcnn: within-site 0.684 -> fold mean 0.608
\newcommand{\xsfDgcnnWithin}{0.684}
\newcommand{\xsfDgcnnFoldMean}{0.608}
\newcommand{\xsfDgcnnAtten}{7.5}  % percentage points, within-site minus fold mean
% pointnet2: within-site 0.622 -> fold mean 0.621
\newcommand{\xsfPtwoWithin}{0.622}
\newcommand{\xsfPtwoFoldMean}{0.621}
\newcommand{\xsfPtwoAtten}{0.0}  % percentage points, within-site minus fold mean

\newcommand{\xsfGoneCellsAbove}{9}  % of 12 cells, folds PH/FC/PR
\newcommand{\xsfGoneNcells}{12}
\newcommand{\xsfGtwoCellsAbove}{5}  % of 12 cells, folds GW/PT/LC
\newcommand{\xsfGtwoNcells}{12}

\newcommand{\xsfNfolds}{6}
\newcommand{\xsfNencoders}{4}
\newcommand{\xsfNruns}{24}  % fresh training runs behind this block


\input{numbers_xsitefix_bin}
%   numbers_errconc.tex  MACHINE-GENERATED, DO NOT HAND-EDIT. Error-concentration
%                        statistics emitted by scripts/paper_reruns/error_concentration.py
%                        from the cached prediction arrays. Replaces the draft's
%                        unsupported "under an independence model..." sentence.
\input{numbers_errconc}
%   numbers_recog.tex    MACHINE-GENERATED, DO NOT HAND-EDIT. Class-size/F1
%                        correlation, matched-size spread and the RASopathy
%                        confusion enrichment, from scripts/paper_reruns/recognizability.py.
%                        RESOLVES the poster-vs-results_analysis rho conflict.
\input{numbers_recog}
%   numbers_perclass_trainseed.tex
%                        MACHINE-GENERATED, DO NOT HAND-EDIT. Added 2026-09-08.
%                        The per-class decomposition of the softmax-vs-prototype
%                        contrast, over the FIVE-TRAINING-SEED sweep. Emitted by
%                        scripts/paper_reruns/emit_perclass.py --family trainseed,
%                        which cross-checks its macro-F1 total against
%                        contrast_tests_trainseed.csv on every run and refuses to
%                        write if the two disagree.
%                        SUPERSEDES the hand-maintained \dxHead* / \synPerClass*
%                        values below and in numbers_perclass.tex, which came
%                        from ONE training run and decomposed a 6.1 pp macro-F1
%                        difference the manuscript no longer reports (it reports
%                        3.6 pp). The old macros are left defined but are no
%                        longer used by the prose.
% ==========================================================================
% numbers_perclass_trainseed.tex — MACHINE-GENERATED. Do not hand-edit.
%   generator : facebench/eval/emit_perclass.py --family trainseed
%   arm       : B{1,2}_corrected_ts{1..5}, five independently
%               trained models; inference seeds averaged down
%               WITHIN each before averaging across.
%   counts    : computed on the MEAN per-class difference. The
%               per-run range is emitted beside every count --
%               quoting one without the other overstates how
%               stable the sorting is.
% ==========================================================================
% ---- B1_corrected (19 classes), 5 training run(s) ----
\newcommand{\tsSynHeadClassesUp}{11}
\newcommand{\tsSynHeadClassesDown}{6}
\newcommand{\tsSynHeadClassesFlat}{2}
\newcommand{\tsSynHeadClassesDead}{0}
\newcommand{\tsSynHeadFoneFromLarge}{1.1}
\newcommand{\tsSynHeadFoneFromSmall}{0.0}
\newcommand{\tsSynNLarge}{19}
\newcommand{\tsSynNSmall}{0}
\newcommand{\tsSynPerClassFoneMad}{0.028}
\newcommand{\tsSynPerClassAucMad}{0.020}
\newcommand{\tsSynPerClassRatio}{1.4}
\newcommand{\tsSynHeadClassesUpRange}{9--12}
\newcommand{\tsSynHeadClassesDownRange}{4--8}
\newcommand{\tsSynHeadClassesFlatRange}{1--3}
\newcommand{\tsSynHeadClassesDeadRange}{0--0}
% ---- B2_corrected (33 classes), 5 training run(s) ----
\newcommand{\tsDxHeadClassesUp}{19}
\newcommand{\tsDxHeadClassesDown}{9}
\newcommand{\tsDxHeadClassesFlat}{3}
\newcommand{\tsDxHeadClassesDead}{2}
\newcommand{\tsDxHeadFoneFromLarge}{2.9}
\newcommand{\tsDxHeadFoneFromSmall}{0.7}
\newcommand{\tsDxNLarge}{25}
\newcommand{\tsDxNSmall}{8}
\newcommand{\tsDxPerClassFoneMad}{0.083}
\newcommand{\tsDxPerClassAucMad}{0.033}
\newcommand{\tsDxPerClassRatio}{2.6}
\newcommand{\tsDxHeadClassesUpRange}{14--17}
\newcommand{\tsDxHeadClassesDownRange}{9--14}
\newcommand{\tsDxHeadClassesFlatRange}{0--2}
\newcommand{\tsDxHeadClassesDeadRange}{4--6}



% ── Superseded RASopathy enrichment (named so the prose can cite what it is
%    retracting without a bare literal). Both values are regenerable with
%    scripts/paper_reruns/recognizability.py --family legacy, which reproduces
%    them exactly; they come from ONE training checkpoint. The five-model
%    values that replace them are \rasWithinFrac / \rasEnrich in
%    numbers_recog.tex, and the per-model range is \rasWithinFracLo--Hi.
\newcommand{\rasWithinFracLegacy}{27\%}   % one checkpoint, 5 inference seeds
\newcommand{\rasEnrichLegacy}{3}          % 2.8x, rounded as published

% ── Cohort ──────────────────────────────────────────────────────────────────
% F2 RESOLVED 2026-08-10. Every count below was verified directly against the
% manifests on disk (row counts and groupby, not transcribed from any document).
%
% THE BUG THAT WAS HERE. The Abstract and Introduction paired \NScans{26,673}
% with \NSubjects{18,076}. Those are TWO DIFFERENT CORPORA and neither number
% described the other:
%   26,673 = the five non-TJ0 datasets counted WHOLE + FB-TJ0 counted at its
%            19-class-FILTERED size. A hybrid. It is not the size of any corpus
%            this paper uses, trains on, or reports — writer's flag F2.
%   18,076 = subjects of combined_manifest.csv, whose scan count is 24,973.
%            A real number, but of the COMBINED BINARY TASK corpus, not of the
%            processed corpus the sentence claimed to describe.
%
% THE CONSISTENT PAIR, verified per dataset (scans / subjects):
%   FB-5A 3,099/3,099 · FB-56 3,839/3,838 · FB-TJ0 13,307/5,103
%   FB-TK0 686/686 · FB-VWP 2,454/2,454 · FB-TX4 3,605/3,605
%   TOTAL 26,990 / 18,785   <- \NScansProc / \NSubjectsProc
% Sources: ofc_manifest.csv (groupby dataset; study_id unique within dataset),
% syndrome_manifest.csv (fbid), controls_manifest.csv (one scan per subject,
% cross-checked against combined_manifest.csv's subject_id assignment).
%
% USE \NScansProc/\NSubjectsProc for the corpus the benchmark DRAWS ON, and
% \combNScansTask/\NSubjects for the combined binary TASK. Never mix the two in
% one sentence — that is exactly the error being fixed here.
\newcommand{\NSubjectsProc}{18{,}785}     % subjects of the 26,990-scan processed corpus
\newcommand{\NSubjects}{18{,}076}         % subjects of the 24,973-scan COMBINED BINARY task
\newcommand{\NScansRound}{${\sim}$27{,}000}  % 26,990 — still correct
\newcommand{\NDatasets}{six}
% \NVendors DELETED 2026-08-10 (writer's flag F3, resolved). "Four scanner
% vendors" had no source: dataset_usage_review.md records 3dMD for FB-5A,
% FB-56 and FB-TJ0 and identifies no scanner for FB-TK0/VWP/TX4. The paired
% claim that "only FB-TJ0 carries colour" was also WRONG — the only textured
% dataset in the review is FB-1A (.obj + .mtl), which is EXCLUDED from this
% paper entirely. Both call sites now state what the metadata supports and
% assign no vendor count. Left undefined so a stale draft fails loudly.
\newcommand{\dsFiveA}{3{,}099}      % FB-5A   | ofc_manifest.csv
\newcommand{\dsFiveSix}{3{,}839}    % FB-56   | ofc_manifest.csv
\newcommand{\dsTJZero}{12{,}990}    % FB-TJ0 at its 19-class-filtered size | syndrome_manifest_b1.csv
\newcommand{\dsTKZero}{686}         % FB-TK0  | controls_manifest.csv
\newcommand{\dsVWP}{2{,}454}        % FB-VWP  | controls_manifest.csv
\newcommand{\dsTXFour}{3{,}605}     % FB-TX4  | controls_manifest.csv

% Per-dataset SUBJECT counts, same sources as the scan counts above. Verified
% 2026-08-10; these are what Table 1 and Fig. 1(a) both read, so the table and
% the figure cannot disagree. FB-TJ0 is the one cohort where scans and subjects
% differ materially (13,307 / 5,103) — that gap is the whole reason the unit of
% analysis matters, so it must never be quoted from memory.
\newcommand{\dsFiveASubj}{3{,}099}
\newcommand{\dsFiveSixSubj}{3{,}838}
\newcommand{\dsTJZeroSubj}{5{,}103}
\newcommand{\dsTKZeroSubj}{686}
\newcommand{\dsVWPSubj}{2{,}454}
\newcommand{\dsTXFourSubj}{3{,}605}

% ── Registration attenuation (Fig. 1d) ──────────────────────────────────────
% Mean per-vertex deviation from the normative template, before and after a
% regularised non-rigid morph of a real 22q11.2 scan toward the real control
% template the pipeline selected. Source: figures/fig_reg_meta.json, produced by
% abstracts/make_registration_fig.py. ILLUSTRATIVE of what dense/non-rigid
% registration does — our own pipeline performs NO registration at all, and the
% caption must keep saying so.
\newcommand{\regDevRaw}{3.6}
\newcommand{\regDevPost}{1.9}
\newcommand{\dsOneAexcl}{\ver{63}}        % FB-1A  — in inventory, in no manifest
\newcommand{\dsVPWexcl}{\ver{26}}         % FB-VPW — in inventory, in no manifest

% F1 RESOLVED 2026-08-10. Was "eight" until writer corrected it 2026-08-04;
% re-verified directly against the manifests, so \ver{} is dropped:
%   syndrome_manifest_b1.csv  12,990 rows / 4,955 subjects, max 34, mean 2.62
%   syndrome_manifest.csv     13,307 rows / 5,103 subjects, max 34, mean 2.61
%   syndrome_clinical_manifest.csv                          max 26
% "Eight" was where the tail thins (37 subjects have exactly 8), not where it
% ends — the old value was not right under any cohort definition.
\newcommand{\maxScansPerSubj}{34}
\newcommand{\medScansPerSubj}{2.6}   % FB-TJ0 mean scans/subject (19-class manifest)

% ── Per-task corpus sizes (Methods: Cohort, Label provenance) ───────────────
% NOTE the three totals are NOT the same corpus and must not be conflated:
%   processed scans across six datasets   = 26,990   (\NScansProc)
%   19-class syndrome-category task       = 12,990   (\synNScans)
%   33-class clinical-diagnosis task      =  3,623   (\dxNScans)
%   combined binary task                  = 24,973   (\combNScansTask)
% \NScans (26,673) HAS BEEN DELETED — see the F2 note in the Cohort block above.
% It was a hybrid matching no corpus and it was the Abstract's headline. Every
% call site now uses \NScansProc. The macro is deliberately left undefined so a
% stale draft quoting it fails at build time rather than shipping the hybrid.
% Row counts verified on disk 2026-08-10: 26,990 = 6,938 (ofc) + 13,307
% (syndrome) + 6,745 (controls); 24,973 = combined_manifest.csv.
\newcommand{\NScansProc}{26{,}990}
\newcommand{\synNScans}{12{,}990}
\newcommand{\synNSubj}{4{,}955}
\newcommand{\dxNScans}{3{,}623}
\newcommand{\dxNSubj}{1{,}241}
\newcommand{\combNScansTask}{24{,}973}
\newcommand{\dsTJZeroAll}{13{,}307}  % FB-TJ0 | syndrome_manifest.csv rows (before the 19-class floor)
\newcommand{\dsTJZeroBin}{11{,}290}  % FB-TJ0 | combined_manifest.csv, dataset==FB-TJ0 (Uncategorized dropped)

% ── Preprocessing attrition (scans in a manifest with no .npy on disk) ──────
\newcommand{\attritSyn}{\ver{12}}
\newcommand{\attritDx}{\ver{5}}
\newcommand{\attritComb}{\ver{361}}

% ── Class inventory ─────────────────────────────────────────────────────────
\newcommand{\synFloorScans}{175}          % min scans/class, 19-class scheme
\newcommand{\dxFloorScans}{50}            % min scans/class, 33-class scheme
\newcommand{\synClassSubjMax}{\ver{978}}  % Control
\newcommand{\synClassSubjMin}{\ver{51}}   % DNA Repair
\newcommand{\dxClassSubjMax}{\ver{123}}   % Down Syndrome
\newcommand{\dxClassSubjMin}{\ver{8}}     % Bohring-Opitz / MPS IIIB
\newcommand{\uncatSubj}{\ver{709}}        % "Uncategorized", 2nd largest 19-class class
\newcommand{\controlSubj}{\ver{978}}      % "Control" category, 19-class scheme
\newcommand{\dysmorphSubj}{\ver{58}}      % "Facial Dysmorphia, possibly genetic"
% RECOMPUTED 2026-09-14 (thread facebase3d-leakpct-2026-09-14): was 86\%, which
% reproduces from nothing. Recomputed from rerun_scanlevel.legacy_split() on
% syndrome_clinical_manifest.csv, keyed on fbid:
%   test SUBJECTS also in train : 364/427 = 0.8525 -> 85\%   <- the sentence's claim
%   test SCANS whose subj in train: 452/530 = 0.8528 -> 85\%  (= \leakDxLeakFrac)
% Neither rounds to 86. Variants checked and rejected: train|val gives 90\%,
% restricting to subjects with repeats gives 92\%. 85\% is correct on BOTH
% readings, so the scans/subjects ambiguity in the prose is harmless here.
\newcommand{\leakTestSeen}{\ver{85\%}}    % 33-cls test subjects also seen in training, scan-level split

% ── Test-set sizes (confirmed by the re-run; coder reply 2026-08-07) ────────
% These matched the published tables exactly, which is what establishes that the
% re-run used the same splits, label maps and metric definitions as the original
% evaluation. The three deterministic encoders also reproduced their published
% A1 accuracies to three decimals (GeomMLP 0.5656 / PointNet 0.7473 / DGCNN 0.7677).
\newcommand{\combNtest}{3{,}639}
\newcommand{\dxNTestSubj}{169}
\newcommand{\synNTestSubj}{733}
\newcommand{\dxNTestScans}{475}
\newcommand{\synNTestScans}{1{,}840}

% ── Chance baselines ────────────────────────────────────────────────────────
\newcommand{\chanceSyn}{5.3\%}            % 19-class
\newcommand{\chanceDx}{3.0\%}             % 33-class
\newcommand{\chanceDxFilt}{4.2\%}         % 24-class filtered

% ============================================================================
% PROSE FORMATTING OF EMITTED INTERVALS
% ============================================================================
% The emitter writes intervals as "[lo, hi]" in raw units. The prose reads them
% inline in a parenthetical, where a bare bracket pair is ambiguous, so they are
% re-rendered here as "95% CI lo to hi" in percentage points. VALUES ARE
% UNCHANGED — only the typography. Each line states the emitted value it wraps
% so the pair can be checked by eye against numbers_reruns.tex.
\renewcommand{\aggGainSynCI}{95\% CI $+3.7$ to $+7.2$}    % emitted [3.7, 7.2]
\renewcommand{\aggGainDxCI}{95\% CI $+5.2$ to $+13.1$}    % emitted [5.2, 13.1]

% ── Aggregation deltas the emitter did not cover ────────────────────────────
% The prototype-rule and kNN aggregation contrasts, from the same
% contrast_tests.csv rows (C2_agg_proto / C2_agg_knn11 [pointnet2]), formatted to
% match the two above. Accuracy unless the name says otherwise.
%
% SUPERSEDED, for the record: \aggGainDx was 12 and \aggGainSyn was 4, both from
% single unseeded runs, and the 19- vs 33-class contrast they supported does not
% exist — the effect is comparable at both granularities and, for the prototype
% rule, larger at 19 classes. The old 19-class "null" was an unseeded-FPS
% artefact, not a finding. Do not reinstate it.
\newcommand{\aggGainDxPr}{8.9}             % pp, 33-class, prototype
\newcommand{\aggGainDxPrCI}{95\% CI $+5.1$ to $+13.1$}
\newcommand{\aggGainSynPr}{9.1}            % pp, 19-class, prototype
\newcommand{\aggGainSynPrCI}{95\% CI $+7.0$ to $+11.1$}
% the one genuine null left: kNN at 19 classes gains no top-1 accuracy, though
% its ranking still improves — the vote, not the pooling, is what discards it.
\newcommand{\aggGainSynKnn}{+0.9}          % pp, 19-class, kNN — spans zero
\newcommand{\aggGainSynKnnCI}{95\% CI $-1.2$ to $+2.7$}
\newcommand{\aggGainSynKnnAuc}{+3.1}       % pp, same contrast, macro AUC
\newcommand{\aggGainSynKnnAucCI}{95\% CI $+1.9$ to $+4.1$}

% ── The combined-binary exception (coder reply 2026-08-07, result (b)) ──────
% Aggregation is NOT a uniform gain. On the combined binary task it buys a small
% amount of accuracy and COSTS macro-F1 — the interval excludes zero on the wrong
% side. Any statement that subject aggregation is a free gain must carry this
% scope limit. Source: contrast_tests.csv, C2_agg_softmax[pointnet2] @ C_corrected.
\newcommand{\aggGainCombAcc}{+0.7}         % pp — SURVIVES
\newcommand{\aggGainCombAccCI}{95\% CI $+0.2$ to $+1.3$}
\newcommand{\aggGainCombFone}{$-1.4$}      % pp — REVERSES
\newcommand{\aggGainCombFoneCI}{95\% CI $-2.1$ to $-0.7$}
\newcommand{\aggGainCombAuc}{$-0.2$}       % pp — DISSOLVES
\newcommand{\aggGainCombAucCI}{95\% CI $-0.7$ to $+0.3$}

% ── Head contrast, prototype minus softmax ──────────────────────────────────
% !! SIGN CONVENTION. contrast_tests.csv and the emitted \headDelta* macros read
% SOFTMAX MINUS PROTOTYPE. The prose reads the other way round, because the
% prototype rule is the intervention under test. These macros are therefore the
% emitted values NEGATED, and each states its emitted counterpart. Never mix the
% two families in one sentence.
%
% Encoder-fixed (contrast C1a): SAME frozen checkpoint, same patients, only the
% inference-time decision rule differs. Intervals are bootstrap over test
% SUBJECTS (n=169 at 33 classes), not the inference-seed spread.
\newcommand{\dxHeadAccDiff}{+3.2}                          % $=-\headDeltaDx$
\newcommand{\dxHeadAccCI}{95\% CI $-2.0$ to $+8.6$}        % emitted [-8.6, 2.0]
\newcommand{\dxHeadFoneDiff}{6.1}
\newcommand{\dxHeadFoneCI}{95\% CI $+0.8$ to $+10.9$}
\newcommand{\dxHeadAucDiff}{$-1.6$}
\newcommand{\dxHeadAucCI}{95\% CI $-2.4$ to $+0.7$}
\newcommand{\synHeadAccDiff}{+0.8}                         % $=-\headDeltaSyn$
\newcommand{\synHeadAccCI}{95\% CI $-1.1$ to $+2.6$}       % emitted [-2.6, 1.1]
\newcommand{\synHeadFoneDiff}{+2.3}
\newcommand{\synHeadFoneCI}{95\% CI $-0.2$ to $+4.7$}

% !! WHAT "NO DIFFERENCE" IS ENTITLED TO SAY (coder reply 2026-08-07, result (c)).
% Of the head contrasts, exactly ONE is a tight null: 19-class macro-AUC, where
% the interval is narrow enough to exclude a material difference. Everywhere else
% the interval contains zero but is too wide to rule out an effect that would
% matter clinically. Those two situations must not be written with the same
% sentence. "The heads tie" is licensed only for 19-class AUC; elsewhere the
% honest form is "we could not separate them, and the interval is wide".
\newcommand{\synHeadAucDiff}{+0.1}                         % $=-\headDeltaSynAuc$
\newcommand{\synHeadAucCI}{95\% CI $-1.0$ to $+1.2$}       % the one tight null

% As-published arm (contrast C1b): the prototype numbers come from a SEPARATE
% CE-trained checkpoint, so this comparison confounds head choice with training
% run. Kept only to show that the confound is what produced the withdrawn claim —
% at 33 classes it "refutes the tie" while the encoder-fixed contrast on the same
% metric does not. Never quote C1b as a head effect.
\newcommand{\dxHeadAccDiffPub}{+7.0}                       % $=-\headDeltaDxPub$
\newcommand{\dxHeadAccCIPub}{95\% CI $+1.2$ to $+13.1$}

% Per-class decomposition of the 33-class macro-F1 difference, computed from the
% per-seed subject-level prediction files (results/paper_reruns/predictions/,
% B2_corrected__pointnet2__{softmax,proto}__subject__seed0-4.npz), five-seed mean
% per class. Checked specifically to rule out the near-singleton-class artefact.
\newcommand{\dxHeadClassesUp}{14}
\newcommand{\dxHeadClassesDown}{9}
\newcommand{\dxHeadClassesFlat}{10}
\newcommand{\dxHeadFoneFromLarge}{4.8}     % pp of the 6.1, from classes n>2
\newcommand{\dxHeadFoneFromTiny}{1.3}      % pp of the 6.1, from the 8 tiny classes
\newcommand{\dxHeadTinyClasses}{eight}     % classes with <=2 test subjects (11 subjects total)

% ── Leakage audit — BEFORE arm only ─────────────────────────────────────────
% !! SUPERSEDED 2026-08-18. The five macros below are the ORIGINAL single-run,
% unseeded "before" values. They are kept under the standing no-deletion rule and
% because they are what the earlier drafts reported, but NOTHING SHOULD CITE THEM
% ANY MORE. Use \leak*BeforeS / \leak*BeforeCI / \leak*DeltaS from numbers_leak.tex,
% which re-score the ORIGINAL checkpoints on the rebuilt legacy scan-level split
% under the same five seeds and the same subject bootstrap as the corrected arm.
%
% The premise of the old comment — that the re-run pipeline could not reproduce
% these — turned out to be wrong rather than unavoidable. dataset.py can no longer
% BUILD the legacy split (it always groups by subject now), but the split code was
% recovered from commit 1d4ec45 and verified: scoring the old B2 checkpoint on the
% rebuilt split returns 0.5075 at batch 1 and a five-seed mean of 0.512 at the
% pipeline's batch 16, against the 0.509 recorded here.
%
% Original note, kept for the record:
% The "after" values are emitted (\leakSynAfter, \leakDxAfter, \leakDxFoneAfter,
% \leakCombAfter). The "before" values come from a pre-correction, scan-level-split
% run whose checkpoint is not in the frozen set, so the re-run pipeline could not
% reproduce them and they stay provisional single-run numbers. Do not attach a CI
% to a before/after delta computed from a \prov value and a seeded one.
\newcommand{\leakDxBefore}{\prov{0.509}}
\newcommand{\leakDxFoneBefore}{\prov{0.476}}
\newcommand{\leakSynBefore}{\prov{0.343}}
\newcommand{\leakCombBefore}{\prov{0.902}} % binary task is robust to the correction
\newcommand{\leakDxDelta}{\prov{22}}       % pp, 0.509 -> 0.291 seeded

% ── Cross-site design constants ─────────────────────────────────────────────
% VERIFIED BY FRANK 2026-08-08 directly against data/processed/ofc_manifest.csv,
% independently of coder's reply, because these went into the manuscript before
% the experiment returned. Reproduced: 6,938 scans / 6,889 subjects / 11 sites;
% ZERO subjects appear at more than one site; 6 sites clear the >=5-scans-in-every-
% one-of-4-classes bar (PH 1,960 / FC 1,259 / GW 744 / PR 589 / PT 473 / LC 126).
% Multi-dataset folds PH, FC, PR (FB-56 + FB-5A) isolate site from dataset;
% GW, PT, LC are single-dataset. The 5 excluded sites are NOT missing a class
% outright in every case — CO and SF have no Cleft Palate, but NG (4/4/3), TX (1)
% and CL (1) merely fall below the floor. Word it as "below the floor", not "absent".
\newcommand{\xsiteNSites}{eleven}
\newcommand{\xsiteNFolds}{six}
\newcommand{\xsiteNTrainSites}{ten}
\newcommand{\xsiteClassFloor}{five}
\newcommand{\xsiteNScans}{6{,}938}
% CORRECTED 2026-08-10. Was 6,889 — that is `study_id` counted ALONE, which
% merges distinct people: 48 study_ids occur in BOTH FB-56 and FB-5A, which
% are separate studies reusing an identifier scheme. Keyed on the pair
% (dataset, study_id) the cohort has 6,937 subjects for 6,938 scans, i.e.
% exactly one subject with two scans. The bootstrap always used the composite
% key (ofc_manifest_xsite.csv subject_id, 6,937 unique), so no RESULT was
% affected — only this quoted count was. Verified: audit_ofc_split.py.
\newcommand{\xsiteNSubj}{6{,}937}

% ── Cross-site transfer ─────────────────────────────────────────────────────
% RETIRED 2026-08-10. \xsiteAucLo{0.48} and \xsiteAucHi{0.51} lived here behind
% an \unsup marker because NO EXPERIMENT EXISTED BEHIND THEM: they were read off
% runs A1_S1/S2/S3, which were believed to be site splits and are not — no code
% path reads the `site` column, all three trained on the identical full cohort,
% and "S" denoted STRATEGY (weighted sampler / focal loss / both). Their real
% values are emitted as \imbWeighted*, \imbFocal*, \imbBoth*.
%
% The leave-one-site-out experiment has now been RUN (24 training runs, 6 folds
% x 4 encoders, 2026-08-09/10) and its numbers live in numbers_xsite.tex.
% THE OLD MACRO NAMES ARE NOT REUSED AND ARE NOT REDEFINED HERE. Rebinding them
% would silently convert a retracted claim into a measured one, and the two are
% not even the same design. Any \xsiteAucLo/\xsiteAucHi left in a draft is now an
% undefined-control-sequence error, which is the intended behaviour: it fails
% loudly rather than quoting a withdrawn number. Use \xsGone*/\xsGtwo*/\xs<Site><Enc>*.

% ── Recognizability (poster block 6) ────────────────────────────────────────
% CONFLICT RESOLVED 2026-08-10. The poster (rho=+0.56, 28% of variance) and
% results_analysis.md Finding 5 (r=0.256, "weak") were BOTH broadly right and
% were computed on DIFFERENT TASKS — 33-class clinical diagnosis vs the
% 19-class/scan-level setting. Recomputed from the seeded predictions in
% numbers_recog.tex: rho=+0.60 at 33 classes, rho=-0.12 at 19. The old
% single-number macros \rhoSize and \varExpl are deleted rather than rebound,
% because there is no single number — the task dependence IS the finding.
% \fOneSpanLo/\fOneSpanHi are now emitted against an explicitly stated size band.

% ── Inference non-determinism ───────────────────────────────────────────────
% The old \fpsWobble / \fpsRunA / \fpsRunB described two ad-hoc UNSEEDED runs and
% are deliberately retired rather than renamed: they were not a measurement of
% anything. Their replacements are emitted from the seeded five-seed sweep —
% \fpsSeedSD (1.07 pp on the headline row), \fpsSeedLo/\fpsSeedHi, \fpsSeedSpan,
% and \fpsMaxSeedSD (3.57 pp, the worst row anywhere in the sweep: 33-class kNN
% subject-level macro-F1, an 8.4 pp span from the sampling draw alone).
% \fpsWobbleOld is kept ONLY for the sentence in Methods that describes what the
% original diagnosis looked like before the path was seeded.
\newcommand{\fpsWobbleOld}{\prov{2--3}}
\newcommand{\fpsRunA}{\prov{0.396}}
\newcommand{\fpsRunB}{\prov{0.367}}

% ── What the published/re-eval discrepancy is, and is NOT ───────────────────
% RESOLVED 2026-08-10 from coder's PART B diagnostics. Every value below was
% re-verified against the raw CSVs, not taken from the summary prose.
%
% THE DISCREPANCY IS TWO-SIDED, so no biased estimator can explain it. Of the 13
% published PointNet++ point estimates, 8 fall outside the five-seed range —
% 5 above and 3 below — while ALL 13 fall inside the bootstrap subject interval.
% (published_direction_check.csv, 13 rows. NOTE: an earlier note in this file and
% coder's 08-07 message both said "10 of 13"; that count was never itemised and
% is withdrawn. This one names task, metric, published value and source line per
% row and is reproducible from check_published_direction.py.)
%
% RULED OUT AS EXPLANATIONS, both with zero effect on any metric:
%   DEVICE (H100 vs CPU) — accuracy identical to 4 dp on every seed at both
%     levels; largest deviation on ANY metric/seed 0.019 pp, ~50x under the
%     1.07 pp inference-seed SD, and the sign flips across seeds.
%     0-1 of 475 test scans change predicted class. (gpu_cpu_deltas.csv)
%   BLAS THREAD COUNT — metrics identical to six decimal places at 1, 8 and 16
%     threads; 0 argmax flips. (diff_threads_20260809.log)
%
% NOT A BIAS, BUT A LARGE SOURCE OF SPREAD: EVAL BATCH SIZE. The FPS start index
% is drawn per sample per batch, so rebatching genuinely resamples the point
% cloud and batch size is part of the estimator, not a throughput knob. Batch 16
% vs 32 changes the predicted class of 94 and 110 of 475 test scans (20-23%),
% yet the aggregate metrics move only within noise and flip sign across seeds
% (accuracy -0.46 pp, macro-F1 -0.38, macro-AUC +0.09; sd exceeds the mean on
% all three). (batchsize_probe.csv, diff_batchsize_20260809.log)
% ONE TASK ONLY (B2_corrected) — do not write that this generalises.
\newcommand{\pubOutside}{eight}
\newcommand{\pubTotal}{thirteen}
\newcommand{\pubAbove}{five}
\newcommand{\pubBelow}{three}
\newcommand{\devMaxDelta}{0.02}      % pp, max |GPU-CPU| on any metric/seed
\newcommand{\threadDP}{six}          % decimal places of agreement across thread counts
\newcommand{\bsFlipLo}{94}
\newcommand{\bsFlipHi}{110}
\newcommand{\bsFlipN}{475}
\newcommand{\bsFlipPct}{20--23}

% ── Encoder spread on the fine-grained tasks ────────────────────────────────
% max - min over the four encoders, scan level, softmax, from summary_ci.csv.
% Derived by subtraction from values already emitted in numbers_reruns.tex, so
% they are checkable against the table rather than taken on trust:
%   19-class accuracy  geommlp 0.115 -> pointnet2 0.295 = 18.1 pp
%   33-class accuracy  pointnet 0.149 -> pointnet2 0.291 = 14.1 pp
%   19-class macro-AUC geommlp 0.650 -> dgcnn 0.781     = 13.2 pp
%   33-class macro-AUC geommlp 0.768 -> dgcnn 0.853     =  8.5 pp
% NOTE the Discussion previously said "roughly ten percentage points" for this
% spread. That was wrong on accuracy and it mattered: the argument there is that
% architecture matters LESS than the evaluation choices, and understating the
% architecture spread made that contrast look sharper than the data supports.
\newcommand{\encSpreadAccLo}{14}   % pp, 33-class accuracy
\newcommand{\encSpreadAccHi}{18}   % pp, 19-class accuracy
\newcommand{\encSpreadAucLo}{9}    % pp, 33-class macro-AUC
\newcommand{\encSpreadAucHi}{13}   % pp, 19-class macro-AUC

% ── Label/cohort structure of the two screening tasks ───────────────────────
% Emitted by scripts/paper_reruns/audit_combined_controls.py --tex, which reads
% the two manifests directly. Re-run it if either manifest changes.
%
% WHY THESE ARE IN THE PAPER (2026-08-17). The COMBINED screening task
% draws 46% of its negative class from three normative cohorts that contribute
% ZERO positives, so cohort membership alone separates the classes at macro-AUC
% 0.730 without reading a single scan. The CLEFT task has no such structure:
% both its cohorts contain both classes, and the same cohort-only rule sits at
% exactly 0.500. That is why Section 2.2 now leads screening with the cleft
% number and reports the pooled figure as an upper bound.
%
% \combCohortOnlyAuc IS NOT A MODEL RESULT AND MUST NEVER BE SUBTRACTED FROM
% ONE. The encoder receives geometry and never a dataset label; this is the
% score of an oracle reading the manifest column. It bounds what COULD be
% cohort recognition. It does not measure what the model does.
\newcommand{\combNegN}{14{,}683}
\newcommand{\combNegNormN}{6{,}745}
\newcommand{\combNegNormFrac}{46\%}
\newcommand{\combCohortOnlyAuc}{0.730}
\newcommand{\combCohortOnlyAcc}{0.682}
\newcommand{\ofcCohortOnlyAuc}{0.500}

% ── Cross-site attenuation of the BINARY screening task ─────────────────────
% Within-site macro-AUC on A2 (numbers_reruns.tex) minus the mean over the six
% leave-one-site-out folds on the SAME task (numbers_xsite_bin.tex). Derived by
% subtraction from emitted values, so it is checkable against the figure:
%   pointnet2  within 0.849 -> fold mean 0.713 = 13.6 pp
%   dgcnn      within 0.805 -> fold mean 0.726 =  7.9 pp
%   geommlp    within 0.649 -> fold mean 0.596 =  5.2 pp
%   pointnet   within 0.543 -> fold mean 0.521 =  2.2 pp
% ONLY the first two are quoted. The last two are attenuations of a within-site
% value that barely clears chance to begin with, so a small drop there is not
% the same quantity as a small drop from 0.85 and must not be averaged with it.
% A FOLD MEAN IS USED ONLY FOR THIS DELTA, never as a headline: the fold ranges
% are what the prose and Fig. 4 report, because the folds are not exchangeable
% (they differ in size, prevalence and whether dataset moves with site).
\newcommand{\xsbAttenLo}{8}    % pp, DGCNN  (7.9, rounded up so the range is not overstated)
\newcommand{\xsbAttenHi}{14}   % pp, PointNet++ (13.6)
\newcommand{\xsbPtwoFoldMean}{0.713}
\newcommand{\xsbDgcnnFoldMean}{0.726}
% Cells above chance in each fold group, ALL FOUR ENCODERS (the \xsb* block
% emits the per-encoder census and the PointNet++ group ranges, not this split).
% Counted from summary_ci_xsite_bin.csv: CI lower bound > 0.50.
\newcommand{\xsbGoneCellsAbove}{11}   % of 12, folds PH/FC/PR
\newcommand{\xsbGtwoCellsAbove}{6}    % of 12, folds GW/PT/LC
% Range of the per-fold majority-class (Unaffected) rate, min and max over the
% six \xsb<Site>Maj values emitted in numbers_xsite_bin.tex:
%   LC 0.708 (min) ... FC 0.919 (max); PH 0.800, PR 0.778, GW 0.812, PT 0.732
\newcommand{\xsbMajLo}{71\%}
\newcommand{\xsbMajHi}{92\%}

% ── Per-class floor ─────────────────────────────────────────────────────────
\newcommand{\nBelowFloor}{\ver{8}}         % of 33 classes with <20 subjects
\newcommand{\nTinyTest}{\ver{12}}          % of 33 classes with <=3 test subjects
\newcommand{\controlRelativeFrac}{\ver{89\%}}

% ── Literature ──────────────────────────────────────────────────────────────
% CITATION AUDIT 2026-08-10. Every value below was read from the source
% PDF in projects/facebase_3d_encoder/papers/ and quoted with its metric and
% class count. Two errors were found and corrected; both had reached the
% Discussion.
%
% ERROR 1 — MISATTRIBUTION. \litMahdiAcc{87.2%} was credited to Mahdi 2024. It
%   is HALLGRIMSSON et al. 2020's number, appearing in Mahdi's related-work section, and it
%   is a TOP-TEN accuracy, not top-1: "the correct diagnosis was listed among
%   the top ten ranked diagnoses for 87.2% of syndromic subjects" (Hallgrimsson et al. 2020;
%   the identical sentence is paraphrased in Mahdi 2024). It was being compared
%   against our TOP-1. Macro deleted; \litHallgTopTen carries it correctly.
%   Mahdi 2024 reports no comparable 33-class top-1 figure — its results table
%   gives sensitivity / balanced accuracy / specificity for outlier-detection
%   models (the 0.8630 that appears there is a BASELINE SPECIFICITY, unrelated).
%
% ERROR 2 — NON-COMPARABLE DENOMINATOR. Bannister 2022a's headline 71% is
%   "overall top-1 accuracy" on a cohort that INCLUDES unaffected subjects, of
%   whom 92% are correctly classified. Our 33-class task is syndromic-only
%   (verified: no unaffected/control class among the 33 in
%   syndrome_clinical_manifest.csv), so that 71% is not a like-for-like target.
%   The paper's own per-syndrome figure is a MEAN SENSITIVITY OF 43% across 48
%   classes, and that is the comparable quantity.
%
% CONSEQUENCE: the old \gapToRegLo{31} / \gapToRegHi{47} "gap to registered
% shape models" was measured against an inflated denominator and a misattributed
% top-10 number. Both are deleted. The honest comparator is Hallgrimsson et al. 2020's
% 48.8% top-1 over syndromic subjects on 64 classes. THIS MOVES IN OUR FAVOUR,
% so it is stated conservatively and the class-count asymmetry (they use more
% classes, which is harder) is kept in the prose.
\newcommand{\litHallgAcc}{48.8\%}       % top-1, syndromic subjects, 64 syndromes (n>=10)
\newcommand{\litHallgNCls}{64}
\newcommand{\litHallgTopTen}{87.2\%}    % TOP-TEN, same cohort — never compare to our top-1
\newcommand{\litBannisterAAcc}{71\%}      % OVERALL top-1, INCLUDES unaffected — see ERROR 2
\newcommand{\litBannisterASens}{43\%}     % mean sensitivity across 48 syndrome classes
\newcommand{\litBannisterANCls}{48}
\newcommand{\litBannisterCAcc}{40.7\%}    % mean top-1 sensitivity, 3D surface PCA+MLP
\newcommand{\litBannisterCNCls}{43}
% VERIFIED 2026-08-19. Read from the source article, not the pointer note.
%   Bannister et al., "Detecting 3D Syndromic Faces as Outliers using Unsupervised
%   Normalizing Flow Models", Artificial Intelligence in Medicine 134:102425 (2022).
%   PMC10949379. Abstract: "the best model achieving an ROC-AUC of 86.3% on a
%   challenging, age and sex diverse data set." Table 3 gives that cell as the
%   NON-GAUSSIAN density model at 1k vertices, 86.3 (1.3) -- the SD is 1.3 pp, which
%   is what makes "a tie" the honest reading against our \ofcBinAuc{} 0.849.
%   Task: binary syndromic vs. non-syndromic. Cohort: 2,471 unaffected + 1,629
%   syndromic across 262 syndromes, ages 5-40, 3dMD scans via FaceBase.
%   Method is dense-correspondence registration + normalizing-flow density
%   estimation -- NOT a statistical shape model; the table's Registration column and
%   the Section~\ref{sec:lit} prose were corrected to match on the same date.
%   NOT the same number as Mahdi's 0.8630, which is a BASELINE SPECIFICITY (checked
%   against the local PDF, Table 2, PCA+LDA baseline row) -- see ERROR 1 above.
\newcommand{\litBannisterBAuc}{0.863}

% Percentage-formatted wrappers for the literature comparison ONLY. The lit
% values are published as percentages, and mixing "0.414" with "48.8%" in one
% sentence invites a reader to compare across a factor of 100. These are
% PRESENTATION of the emitted macros, not new values — if the emitted value
% changes, these must be re-derived from it by hand, so keep the pairing visible.
\newcommand{\dxAccSubjSMPct}{38.2\%}   % = \dxAccSubjSM{} 0.382
\newcommand{\dxAccSubjPrPct}{41.4\%}   % = \dxAccSubjPr{} 0.414
\newcommand{\combAucPct}{96.0\%}       % = \combAuc{} 0.960
